\section*{第 \ref{ch:g10:signals-and-waves} 章——信号与波}

\begin{solution}{exo:g10:signals-and-waves:1}
$f = 1/T$：(a) $1/\qty{0.020}{s} = \qty{50}{Hz}$；
(b) $1/\qty{4.0e-3}{s} = \qty{250}{Hz}$；
(c) $1/\qty{5.0e-5}{s} = \qty{20}{kHz}$。
\end{solution}

\begin{solution}{exo:g10:signals-and-waves:2}
$T = 1/f$：(a) \qty{20}{ms}；(b) $1/440 = \qty{2.3}{ms}$；
(c) $1/\qty{2.4e9}{Hz} = \qty{0.42}{ns}$。
\end{solution}

\begin{solution}{exo:g10:signals-and-waves:3}
\qty{12}{Hz}：\omterm{def:g10:signals-and-waves:ultrasound}{次声波}。\qty{440}{Hz} 和 \qty{18}{kHz}：可听声（
后者仅年轻人能听到）。\qty{40}{kHz} 和 \qty{5}{MHz}：\omterm{def:g10:signals-and-waves:ultrasound}{超声波}。
\end{solution}

\begin{solution}{exo:g10:signals-and-waves:4}
$T = 5.0 \times \qty{2}{ms} = \qty{10}{ms}$，故 $f = \qty{100}{Hz}$；
\omterm{def:g10:signals-and-waves:amplitude}{振幅}为 $4.0 \times \qty{0.5}{V} = \qty{2.0}{V}$。
\end{solution}

\begin{solution}{exo:g10:signals-and-waves:5}
$d = 340 \times 6.0 = \qty{2040}{m} \approx \qty{2.0}{km}$。光传播
这段距离只需 $2040/(3.00\times10^{8}) \approx \qty{7}{\micro s}$，
比 \qty{6.0}{s} 小一百万倍：闪光标志着雷击发生的时刻。
\end{solution}

\begin{solution}{exo:g10:signals-and-waves:6}
\omterm{def:g10:signals-and-waves:sound}{声音}先向下传播，再\emph{返回}，共走过 $2d$：
$d = \frac{1500 \times 0.60}{2} = \qty{450}{m}$。经过海沟：
$t = 2d/v = 2 \times 1200/1500 = \qty{1.6}{s}$。
\end{solution}

\begin{solution}{exo:g10:signals-and-waves:7}
$d_1 = \frac{3.00\times10^{8} \times 2.4\times10^{-4}}{2} =
\qty{36.0}{km}$，$d_2 = \qty{35.7}{km}$：飞机在
\qty{1.0}{s} 内飞过 \qty{300}{m}，故 $v = \qty{300}{m/s} \approx
\qty{1100}{km/h}$，正在靠近。
\end{solution}

\begin{solution}{exo:g10:signals-and-waves:8}
$T = \qty{20}{mm}/\qty{25}{mm/s} = \qty{0.80}{s}$，故
$f = \qty{1.25}{Hz}$：每分钟 $75$ 次搏动。每分钟 $100$ 次搏动时，
$T = \qty{0.60}{s}$，即 $25 \times 0.60 = \qty{15}{mm}$：心动过速
表现为低于 \qty{15}{mm}。
\end{solution}

\begin{solution}{exo:g10:signals-and-waves:9}
$T = 1/200 = \qty{5.0}{ms}$；两个\omterm{def:g10:signals-and-waves:period}{周期}持续 \qty{10}{ms}，分布在
$10$ 格中：\omterm{def:g10:signals-and-waves:oscilloscope}{时基}为每格 \qty{1}{ms}。波峰须能容纳于
$4$ 格之内：每格\omterm{def:g10:signals-and-waves:oscilloscope}{增益}至少为 $3.0/4 = \qty{0.75}{V}$——
取每格 \qty{1}{V}（波峰高 $3.0$ 格）。
\end{solution}

\begin{solution}{exo:g10:signals-and-waves:10}
$d = v\,t/2$：前壁
$\frac{1540 \times 4.0\times10^{-5}}{2} = \qty{3.1}{cm}$；后壁
$\frac{1540 \times 6.0\times10^{-5}}{2} = \qty{4.6}{cm}$；厚度
约为 \qty{1.5}{cm}。
\end{solution}

\begin{solution}{exo:g10:signals-and-waves:11}
\omterm{def:g10:signals-and-waves:sound}{声音}：$30/340 = \qty{0.088}{s}$。无线电：
$3.0\times10^{6}/(3.00\times10^{8}) = \qty{0.010}{s}$。远处的无线电
听众先听到鼓声，约早 \qty{78}{ms}。
\end{solution}

\begin{solution}{exo:g10:signals-and-waves:12}
蝙蝠在发射 \qty{3.0}{ms} 期间无法接收\omterm{rem:g10:signals-and-waves:honest}{回声}，即
距离在 $d = \frac{340 \times 3.0\times10^{-3}}{2} =
\qty{0.51}{m}$ 以内的\omterm{rem:g10:signals-and-waves:honest}{回声}都无法接收。飞蛾位于 \qty{1.7}{m}：
$t = 2 \times 1.7/340 = \qty{10}{ms}$。飞蛾靠近时，\omterm{rem:g10:signals-and-waves:honest}{回声}返回得越来越
早；脉冲必须在\omterm{rem:g10:signals-and-waves:honest}{回声}到达前结束，因此必须缩短。
\end{solution}

\begin{solution}{exo:g10:signals-and-waves:13}
$d = \frac{340 \times 1.5}{2} = \qty{255}{m}$。在 \qty{155}{m} 处：
$t = 2 \times 155/340 = \qty{0.91}{s}$。\omterm{rem:g10:signals-and-waves:honest}{回声}与拍手声
在 $t < \qty{0.1}{s}$，即距离
$d < \frac{340 \times 0.1}{2} = \qty{17}{m}$ 以内时混在一起。
\end{solution}

\begin{solution}{exo:g10:signals-and-waves:14}
\emph{1.} $T = 3.0 \times \qty{5}{ms} = \qty{15}{ms}$，
$f = 1/0.015 \approx \qty{67}{Hz}$。
\emph{2.} $\qty{15}{ms}/\qty{2}{ms} = 7.5$ 格。
\emph{3.} 屏幕依次显示 \qty{50}{ms} 和 \qty{20}{ms}：完整\omterm{def:g10:signals-and-waves:period}{周期} $3$ 个
（$50/15 = 3.3$），另有 $1$ 个（$20/15 = 1.3$）。
\end{solution}

\begin{solution}{exo:g10:signals-and-waves:15}
\emph{1.} $t = 2 \times 1.2\times10^{6}/(3.00\times10^{8}) =
\qty{8.0}{ms}$。
\emph{2.} $2 \times 1.2\times10^{6}/(2.0\times10^{8}) = \qty{12}{ms}$。
\emph{3.} \omterm{def:g10:refraction:fiber}{光纤}并非在两台机器之间笔直延伸，途中的每个路由器和服务器
都会增加处理延迟。
\end{solution}

\begin{solution}{pb:g10:signals-and-waves:1}
\textbf{1.} $d = 340 \times 9.0 = \qty{3060}{m} \approx \qty{3.1}{km}$。

\textbf{2.} $3060/(3.00\times10^{8}) \approx \qty{10}{\micro s}$，
比 \qty{9.0}{s} 小一百万倍：闪光标志着
雷击发生的时刻。

\textbf{3.} 在 \qty{3}{s} 内\omterm{def:g10:signals-and-waves:sound}{声音}传播 $340 \times 3 = \qty{1020}{m}
\approx \qty{1}{km}$：秒数除以三得到千米数。

\textbf{4.} $d = 340 \times 6.0 = \qty{2040}{m}$。暴风雨
在 \qty{300}{s} 内逼近了 \qty{1020}{m}：$v = \qty{3.4}{m/s} \approx
\qty{12}{km/h}$，正朝小船而来。

\textbf{5.} 近端的\omterm{def:g10:signals-and-waves:sound}{声音}经 $2000/340 = \qty{5.9}{s}$ 后听到，远端
的\omterm{def:g10:signals-and-waves:sound}{声音}经 $5000/340 = \qty{14.7}{s}$ 后听到：隆隆声持续约
\qty{8.8}{s}。

\textbf{6.} $d = 1500 \times 0.90/2 = \qty{675}{m}$。

\textbf{7.} $d = 1500 \times 0.20/2 = \qty{150}{m}$。

\textbf{8.} 两个障碍物：位于 $1500 \times 0.16/2 =
\qty{120}{m}$ 处的鱼群，以及 \qty{150}{m} 处的海底。鱼群
游在海底上方 \qty{30}{m} 处。

\textbf{9.} \omterm{rem:g10:signals-and-waves:honest}{回声}必须在下一次脉冲之前返回：
$d_{\max} = 1500 \times 0.50/2 = \qty{375}{m}$。若水更深，
\omterm{rem:g10:signals-and-waves:honest}{回声}会在下一次脉冲\emph{之后}才到达，并被归到它名下：
显示器便显示出一个虚假的、浅得多的海底。

\textbf{10.} 在空气中，$340 \times 0.60/2 = \qty{102}{m}$，而不是
\qty{450}{m}。只有知道载波及其在所经介质中的速率，
延迟时间才能换算为距离。

\textbf{11.} $f = 1/0.86 = \qty{1.16}{Hz}$，即 $1.16 \times 60
\approx 70$ 次/分。

\textbf{12.} $25 \times 0.86 = \qty{21.5}{mm}$。

\textbf{13.} $T = 15/25 = \qty{0.60}{s}$：$100$ 次/分。

\textbf{14.} $T = 60/50 = \qty{1.2}{s}$；间距为
$25 \times 1.2 = \qty{30}{mm}$。

\textbf{15.} 否：\omterm{def:g10:signals-and-waves:period}{周期}会随着用力、呼吸和压力逐拍变化，
医生读取的是\omterm{def:g10:signals-and-waves:period}{周期}（即心率）、其规律性，以及每个\omterm{def:g10:signals-and-waves:period}{周期}的波形——
诊断的依据就在
\omterm{def:g10:signals-and-waves:signal}{信号}之中。

\textbf{16.} 无线电：$5.0\times10^{4}/(3.00\times10^{8}) =
\qty{0.17}{ms}$——对人的感官而言是瞬时的。\omterm{def:g10:signals-and-waves:sound}{声音}：
$50\,000/340 \approx \qty{147}{s}$，约两分半钟。

\textbf{17.} $t = 2.02\times10^{7}/(3.00\times10^{8}) = \qty{67}{ms}$。

\textbf{18.} $c \times \qty{1.0}{\micro s} = \qty{300}{m}$。对于
\qty{3.0}{m} 的定位：$3.0/(3.00\times10^{8}) = \qty{10}{ns}$——
这正是 GPS \omterm{def:g10:universal-gravitation:satellite}{卫星}携带原子钟的原因。

\textbf{19.} \omterm{def:g10:universal-gravitation:satellite}{卫星}处于真空中，而\omterm{def:g10:signals-and-waves:sound}{声音}需要介质：
\omterm{def:g10:universal-gravitation:satellite}{卫星}根本发不出任何\omterm{def:g10:signals-and-waves:signal}{信号}。（即便假定
全程都有空气，以 \qty{340}{m/s} 的速率传播 \qty{20200}{km} 也需超过 \qty{16}{h}——
即得到的定位已过时数小时。）

\textbf{20.} \qty{340}{m/s}（空气中的\omterm{def:g10:signals-and-waves:sound}{声音}）、\qty{1500}{m/s}（海水中的\omterm{def:g10:signals-and-waves:sound}{声音}）
和 $c = \qty{3.00e8}{m/s}$（无线电波和光）；规律是
$d = v\,t$，\omterm{rem:g10:signals-and-waves:honest}{回声}时取一半。习惯是：知道你的载波，知道它
在介质中的速率，测量延迟——延迟不过是
伪装了的距离。
\end{solution}
